\documentclass[11pt,a4paper]{article}
\usepackage[T1]{fontenc}
\usepackage[utf8]{inputenc}
\usepackage{lmodern,amsmath,amssymb}
\usepackage[margin=25mm]{geometry}
\usepackage[hidelinks]{hyperref}
\hypersetup{pdftitle={Cut-point consistency and an action remainder},pdfauthor={navstokgap research working notes}}
\newcommand{\tightlist}{\setlength{\itemsep}{0pt}\setlength{\parskip}{0pt}}
\setlength{\emergencystretch}{3em}
\setcounter{secnumdepth}{0}
\begin{document}
\hypertarget{a-surviving-action-defect-and-consistency-under-inserting-cuts}{%
\section{A surviving action defect and consistency under inserting
cuts}\label{a-surviving-action-defect-and-consistency-under-inserting-cuts}}

A03, reviewed with B12, 2026-09-07. Two exact tests separate refinement
of one experiment from a changing family of fluctuation preparations.
Constant-force chord errors vanish on every mesh tending to zero. Within
the scalar Gaussian bridge family, exact consistency under retaining old
nodes fixes the fluctuation parameter. These statements locate the next
selection obligation at the refinement law.

\hypertarget{arbitrary-partitions-of-the-constant-force-trajectory}{%
\subsection{1. Arbitrary partitions of the constant-force
trajectory}\label{arbitrary-partitions-of-the-constant-force-trajectory}}

Fix \(m,F,T>0\), \(x(t)=v_0t\) with \(v_0>0\), \(y(t)=Ft^2/(2m)\), and
the Lagrangian \(L=m(\dot x^2+\dot y^2)/2+Fy\). Let
\(\pi=(0=t_0<\cdots<t_N=T)\), \(\tau_j=t_j-t_{j-1}\), and let \(q_\pi\)
linearly interpolate the exact trajectory at these times. Both global
endpoints and all sampled positions are fixed in this comparison.

On interval \(j\), with \(s=t-t_{j-1}\), the interpolation error is
\(\eta_j(s)=Fs(\tau_j-s)/(2m)\). The classical first variation vanishes
on that interval. Adding the exact local quadratic action differences
gives

\[D_\pi:=S[q_\pi]-S[q_{\rm cl}]
=\frac{F^2}{24m}\sum_j\tau_j^3
\le\frac{F^2T}{24m}|\pi|^2\longrightarrow0.\]

The unsigned chord--arc lens areas similarly add to
\(A_\pi=v_0F\sum_j\tau_j^3/(12m)\), so \(D_\pi=F A_\pi/(2v_0)\). The
existing constant-force calculation applies after subtracting the local
linear tangent; its slope cancels from the lens difference.

For fixed \(N\), convexity yields \(\sum_j\tau_j^3\ge T^3/N^2\),
attained by equal steps. Splitting one interval into positive lengths
\(a,b\) reduces the action defect by

\[\frac{F^2}{24m}[(a+b)^3-a^3-b^3]
=\frac{F^2}{8m}ab(a+b)>0.\]

Thus nonuniform cuts cannot defeat convergence when the maximum step
tends to zero. This calculation concerns sampled chords; Newton's
impulsive polygon is a separate approximation scheme requiring its own
error estimate.

\hypertarget{exact-consistency-of-gaussian-node-laws}{%
\subsection{2. Exact consistency of Gaussian node
laws}\label{exact-consistency-of-gaussian-node-laws}}

Fix the free-line experiment on \([0,T]\), mass \(m>0\) and fixed
endpoints. For each partition with at least one internal node, let the
centered node vector be Gaussian with covariance

\[\operatorname{Cov}_{\pi}(\eta_i,\eta_j)
=\frac{\kappa_\pi}{m}G(t_i,t_j),\qquad
G(s,t)=\min(s,t)-st/T,\quad \kappa_\pi\ge0.\]

The zero value denotes a point mass at the classical straight path.
Exact refinement consistency means that when \(\pi'\supset\pi\),
deleting the new coordinates from the finer law gives precisely the
coarser law, with the same physical times, positions and mass.

\textbf{Proposition.} Within this family, consistency forces
\(\kappa_{\pi'}=\kappa_\pi\) for every refinement retaining an interior
node. Conversely a fixed \(\kappa\) gives consistent Gaussian node laws.

\textbf{Proof.} At any retained interior time \(s\), equality of the
one-node marginals implies
\(\kappa_{\pi'}s(T-s)/(mT)=\kappa_\pi s(T-s)/(mT)\). The factor
multiplying \(\kappa\) is positive. Conversely the restriction of a
centered Gaussian vector has the corresponding covariance submatrix,
which here is exactly the coarse covariance. On the directed family of
all partitions, any two nontrivial partitions have a common refinement,
so their parameters agree. The endpoint-only partition has no observable
variance and places no extra constraint. This completes the proof.

For nested \(\pi_N\) retaining an interior point \(s\), M05's
finite-defect scaling \((N-1)\kappa_N\to\ell>0\) makes that retained
variance approach zero. It therefore describes changing preparations,
rather than exact marginals of one positive-width bridge experiment. It
remains a valid triangular-array limit, with the action defect proved in
C014.

\hypertarget{what-survives-at-fixed-preparation}{%
\subsection{3. What survives at fixed
preparation?}\label{what-survives-at-fixed-preparation}}

Let \(d=N-1\). Under the fixed-\(\kappa>0\) bridge, the polygon's free
action excess satisfies \(2D_\pi/\kappa\sim\chi_d^2\), independent of
the actual positive step lengths (C014's finite law). Consequently

\[\mathbb E D_\pi=d\kappa/2,\qquad
\operatorname{Var}D_\pi=d\kappa^2/2,\]

and \(D_\pi\to\infty\) in probability as \(d\to\infty\). Indeed,
Chebyshev gives \(\Pr\{D_\pi<d\kappa/4\}\le8/d\). The normalized
estimator

\[\widehat\kappa_\pi=\frac{2D_\pi}{d}\]

has mean \(\kappa\) and variance \(2\kappa^2/d\), hence converges to
\(\kappa\) in mean square, even on nonuniform partitions. It estimates
the supplied bridge parameter. It is not an additive accumulated action;
division by the number of internal nodes is essential. For \(\kappa=0\)
every defect vanishes.

Thus exact consistency permits a fixed positive action parameter and a
consistent estimator of it, while the unnormalized kinetic action
diverges. It also permits the deterministic zero family. Selection of a
positive scale still requires a physical premise excluding that family
and fixing the scale.

\hypertarget{one-inserted-cut-and-the-fluctuation-it-introduces}{%
\subsection{4. One inserted cut and the fluctuation it
introduces}\label{one-inserted-cut-and-the-fluctuation-it-introduces}}

On a coarse interval of length \(a+b\), fix endpoint positions \(x,z\)
and insert \(y\) after duration \(a\), with \(a,b>0\). Set
\(\bar y=(bx+az)/(a+b)\) and \(\zeta=y-\bar y\). Completing the square
gives

\[\frac m2\left[\frac{(y-x)^2}{a}+\frac{(z-y)^2}{b}
-\frac{(z-x)^2}{a+b}\right]
=\frac{m(a+b)}{2ab}\zeta^2.\]

At fixed bridge parameter \(\kappa\), conditional Gaussian variance is
\(\operatorname{Var}(\zeta\mid x,z)=\kappa ab/[m(a+b)]\). The expected
extra action from inserting one node is therefore \(\kappa/2\),
independent of the interval lengths. For deterministic interpolation
\(\zeta=0\) it is zero. The kinetic splitting identity is nonnegative,
whereas section 1's total constant-force chord error decreases under
refinement: the latter includes the potential term and moves the
inserted node onto the accelerated classical curve.

This is a local, testable formulation of the missing premise: does
inserting a physical cut introduce a nonzero conditional fluctuation,
and what law sets its variance? Supplying the Gaussian rule answers the
first question by assumption. A classical derivation would have to
produce that rule or a different consistent fluctuation mechanism from
independent physics.

\hypertarget{next-test-and-prior-art-gate}{%
\subsection{5. Next test and prior-art
gate}\label{next-test-and-prior-art-gate}}

The \href{../references/batches/B12.md}{B12 audit} classifies C027--C029
as elementary consequences. Pitman and Yor, \emph{A guide to Brownian
motion and related stochastic processes} (2018),
\href{https://arxiv.org/abs/1802.09679v1}{arXiv:1802.09679v1}, printed
pp.~6, 10, 12--13, supply the Gaussian bridge and restriction framework.
The \href{../docs/batches/B12/PitmanYor2018.md}{source companion}
records coverage. The constant-force identity extends the
\href{../notes/principia-constant-force-action.md}{existing
calculation}. The action law and moments in section 3 are proved in
\href{../papers/regulator-limits.tex}{the regulator-limits paper}. The
combined selection test is the research use of these established
ingredients.

Next distinguish coordinate sampling from a physically executed
intervention at a cut. A candidate intervention law must state its
effect on energy, momentum, conditional variance and coarse observables.
Test consistency first, then positivity, scale universality and
compatibility with finite speed. The Gaussian conditional law is a
reference model for that audit, not yet a model with a hard physical
speed ceiling.
\end{document}
